% Compile: tectonic main.tex (or xelatex main.tex twice). Output: main.pdf.
\documentclass[UTF8,fontset=windows,a4paper,zihao=-4]{ctexart}
\usepackage[left=23mm,right=23mm,top=24mm,bottom=24mm,headheight=15pt]{geometry}
\usepackage{amsmath,amssymb,booktabs,array,longtable,graphicx,xcolor,tikz,fancyhdr}
\usetikzlibrary{arrows.meta,positioning,calc,fit,backgrounds}
\usepackage{xurl}
\usepackage[unicode,colorlinks=true,linkcolor=teal!60!black,urlcolor=teal!60!black]{hyperref}
\hypersetup{pdftitle={GESTURE PLAY 系统设计报告},pdfsubject={设计原理、功能框图、软硬件划分与接口设计}}
\definecolor{ink}{HTML}{163247}
\definecolor{accent}{HTML}{087F8C}
\definecolor{pale}{HTML}{EAF4F5}
\setlength{\parskip}{4pt}\linespread{1.18}
\setlength{\emergencystretch}{3em}\setcounter{tocdepth}{2}
\ctexset{section={format=\Large\bfseries\color{ink}},subsection={format=\large\bfseries\color{ink}}}
\pagestyle{fancy}\fancyhf{}
\fancyhead[L]{\small GESTURE PLAY}\fancyhead[R]{\small 初级组 · 系统设计}
\fancyfoot[C]{\small\thepage}\renewcommand{\headrulewidth}{0.3pt}
\tikzset{box/.style={draw=accent,fill=pale,rounded corners=2pt,align=center,minimum height=12mm,text width=30mm,font=\small},soft/.style={box,fill=white,draw=ink!60},arr/.style={-{Stealth[length=2mm]},thick,draw=ink},ctl/.style={arr,dashed},lab/.style={font=\footnotesize,fill=white,inner sep=2pt,align=center}}
\begin{document}
\hypersetup{pageanchor=false}
\begin{titlepage}\color{ink}\vspace*{15mm}
{\large FPGA 创新设计赛道 · 自主选题 · 初级组\par}\vspace{15mm}
{\fontsize{32}{40}\selectfont\bfseries GESTURE PLAY\par}\vspace{3mm}
{\fontsize{25}{34}\selectfont\bfseries 系统设计报告\par}\vspace{8mm}
{\Large 设计原理与功能框图\par}\vspace{3mm}
{\Large 软硬件划分依据与接口设计\par}\vspace{10mm}
\noindent{\color{accent}\rule{\linewidth}{1.2pt}}\par\vspace{6mm}
{\large 基于 PYNQ-Z2 的手势交互、图形显示\par 与灵巧手协同系统\par}\vspace{12mm}
\noindent\colorbox{pale}{\parbox{\dimexpr\linewidth-2\fboxsep\relax}{\vspace{5mm}\centering
\textbf{视觉采集 → 手势识别 → 应用决策}\par
\textbf{图形呈现与实体动作反馈}\par\vspace{4mm}
四游戏 · 命令式二维图形处理器 · 画布相册 · Studio\par\vspace{5mm}}}
\vfill\noindent\begin{tabular}{@{}ll}
硬件标识 & BUILD\_ID 0xC0061000\\[4pt]
器件与工具 & xc7z020clg400-1 / Vivado 2026.1\\[4pt]
系统平台 & PYNQ-Z2 / Python 3.10 / PYNQ 3.1\\[4pt]
更新日期 & 2026 年 11 月 4 日
\end{tabular}\end{titlepage}
\pagenumbering{roman}\hypersetup{pageanchor=true}\tableofcontents\clearpage
\pagenumbering{arabic}\section{设计概述与初级组要求对应}
GESTURE PLAY v2.0 是面向手势交互与创作体验的 Zynq 系统。作品以 PYNQ-Z2 为核心，将摄像头视觉输入、自研 Scene2D 二维图形处理器、四款游戏、画布相册、视频与灵巧手协同组织为一个完整应用。用户通过手的位置、握拳和板载按键完成选择与操作，系统给出 HDMI 图形反馈，并通过独立的实体手通道执行姿态与演示动作。

v2.0 已完成从输入采集、手势识别、应用决策到图形呈现的板端运行链路；PC Studio 提供内容创作、部署、测试与状态观测。发布验收结果为 PASS：四游戏 16 条 PC/PYNQ 对照用例、真实 GPU 60 项检查、真实串口协议 18 项检查，并完成设备关闭回执、正常服务与 Studio 重连验证。[E01–E03]

\subsection{作品构成与实现基线}
作品采用 PYNQ-Z2（xc7z020clg400-1），以 Vivado 2026.1 构建 R23b2 硬件，BUILD\_ID 为 0xC0061000。板端使用 Python 3.10 / PYNQ 3.1 和 R09 灵巧手运行时，发布基线于 2026-10-29 封存。

视觉端输出 640×360 双掩膜；显示端以 320×180 逻辑画面驱动 1280×720 / 60 Hz HDMI。四游戏、两段默认视频、手势菜单、画布相册及实体手/像素手演示共享这一底座，Studio 通过统一 SDK 与 Agent 管理应用。

\subsection{与初级组考察重点的对应}
官方自主选题指南将初级组的重点列为逻辑设计、状态机、时序约束及接口协议，并推荐采用 PYNQ 组织 PS 软件与 PL Overlay。本作品以这些基本能力贯穿完整应用：像素流在 PL 中完成规则化处理，软件通过明确接口组合各模块，最终在真实板卡上形成可操作、可观察、可验证的作品。[R01]

{\small\renewcommand{\arraystretch}{1.28}
\begin{longtable}{@{}>{\raggedright\arraybackslash}p{0.24\linewidth}>{\raggedright\arraybackslash}p{0.35\linewidth}>{\raggedright\arraybackslash}p{0.34\linewidth}@{}}
\toprule
\bfseries 考察重点 & \bfseries 本作品的具体实现 & \bfseries 本报告展开内容\\\midrule
\endfirsthead
\toprule
\bfseries 考察重点 & \bfseries 本作品的具体实现 & \bfseries 本报告展开内容\\\midrule
\endhead
\bottomrule\endfoot
逻辑设计 & 双掩膜流水线、命令光栅化、双缓冲显示、独立 OSD & 数据通路、存储结构、流水对齐\\
状态机 & 帧边界启停、三槽缓冲拥有权、显示提交与模式交接 & 状态转移、完成条件、计数不变量\\
时序约束 & 100 MHz 控制域、74.25 MHz 像素域、跨域握手 & 时钟划分、同步器、数据路径约束与实现结果\\
接口协议 & AXI、DMS1、Scene2D ABI、VVGP、灵巧手观察通道 & 格式、握手、校验与回执\\
作品完整性 & 四游戏和创作功能共享运行底座；服务重启与 Studio 重连 & 功能矩阵、完整事务和实板验收\\
\end{longtable}}
\clearpage
\section{设计原理与功能框图}
\subsection{系统总体结构}
\begin{figure}[htbp]\centering\begin{tikzpicture}[node distance=9mm and 10mm]
\node[box] (cam) {OV5640\\摄像头};
\node[box,right=of cam] (vision) {PL 视觉前端\\双掩膜流水线};
\node[box,right=of vision] (ps) {PS / PYNQ\\识别与应用};
\node[box,below=13mm of vision] (gpu) {PL Scene2D\\渲染与扫描};
\node[soft,left=of gpu] (screen) {HDMI\\图形反馈};
\node[soft,below=13mm of ps] (hand) {独立动作进程\\USB 灵巧手};
\node[soft,above=13mm of ps] (pc) {PC Studio\\创作 / 部署 / 测试};
\draw[arr] (cam)--(vision);
\draw[arr] (vision)--node[lab,above=5mm]{DMA / DDR}(ps);
\draw[arr,<->] (ps.south west)--node[lab,above,sloped]{命令 / 回执}(gpu.north east);
\draw[arr] (gpu)--(screen);
\draw[arr] (ps)--node[lab,right]{观察}(hand);
\draw[ctl,<->] (pc)--node[lab,right]{VVGP}(ps);
\end{tikzpicture}
\caption{系统功能框图}\label{fig:1}\end{figure}
系统由视觉输入、应用运行、图形输出、实体动作和开发管理五部分组成。OV5640 图像进入 PL 后完成预处理与双掩膜生成，经 DMA 送入 DDR；PS 从掩膜提取轮廓、质心及手势状态，并将结果分发给应用、OSD 与灵巧手观察通道。游戏和界面在 PS 中生成 Scene，经 DMA 提交到 PL 的 Scene2D 处理器，最后由独立扫描电路输出 HDMI。[E04–E13]

在数据组织上，像素与命令使用批量通路，配置和状态使用 AXI-Lite 寄存器通路。前者关注吞吐、帧边界及缓冲生命周期，后者关注配置一致性、身份和完成回执。这一划分使高速数据处理与灵活的应用状态管理相互配合。

交互闭环由“用户动作—采集识别—应用响应—屏幕或实体反馈—用户继续操作”组成。系统同时建立软件层面的完成反馈：视觉包携带帧身份，图形提交区分接收、渲染和呈现，实体通道在模式交接时提供关闭读回与会话回执。这些反馈使一次交互既有可见结果，也有可追溯的执行记录。

\subsection{视觉流水线与手势识别}
\begin{figure}[htbp]\centering\begin{tikzpicture}[node distance=8mm and 10mm]
\node[box] (rgb) {RGB 图像流\\1280×720};
\node[box,right=of rgb] (area) {2×2 区域处理\\640×360};
\node[box,right=of area] (hsv) {HSV 阈值判定\\原始二值掩膜};
\node[soft,below=10mm of area] (raw) {原始掩膜 FIFO\\数据 / 帧标志};
\node[box,below=10mm of hsv] (filter) {9×9 膨胀\\5×5 高斯滤波};
\node[box,below=10mm of raw] (pair) {双路对齐与配对\\16 bit / 像素};
\node[box,left=of pair] (dma) {DMS1 + DMA\\身份 / 阈值 / CRC};
\node[soft,below=10mm of dma] (cpu) {PS 轮廓与形状判断\\质心 / 手势 / 事件};
\draw[arr] (rgb)--(area);\draw[arr] (area)--(hsv);
\draw[arr] (hsv)--(filter);\draw[arr] (hsv.south west)--(raw.north east);
\draw[arr] (raw)--(pair);\draw[arr] (filter.south)|-(pair.east);
\draw[arr] (pair)--(dma);\draw[arr] (dma)--(cpu);
\node[lab,below=10mm of pair,text width=63mm,anchor=north west,xshift=-15mm] {上下文 FIFO 随帧传递\\epoch / config\_id / frame\_id / 阈值};
\end{tikzpicture}
\caption{视觉流水线与双掩膜配对}\label{fig:2}\end{figure}
PL 视觉前端以 24 bit 图像流为输入，先由 axis\_area2x2 完成 2×2 区域处理，将 1280×720 图像缩减为 640×360。随后进行颜色字节顺序适配和 HSV 肤色判定。六个 HSV 阈值在一帧的首个有效输入处锁存，同帧像素采用相同配置。[E19]

原始二值掩膜分成两支：一支保留原始判定并进入同步 BRAM FIFO；另一支经过 9×9 膨胀及 5×5 高斯滤波。配对模块在两路数据和帧标志都对应时输出一个 16 bit 像素，其中低字节为原始掩膜，高字节为滤波结果。FIFO 为两条处理延迟不同的通路提供对齐空间，控制逻辑共同推进数据与 tuser/tlast 标志。溢出、下溢和帧标志失配均进入可读故障状态。

视觉前端还设置深度为 4 的上下文 FIFO，每项 144 bit，保存 epoch、config\_id、frame\_id 与六个阈值字节。输入首像素入队、输出帧尾出队，使像素数据与产生它的配置一同穿过流水线。这样，PS 收到的数据不仅包含“哪里是手”，也明确标识“哪一帧、哪一组阈值、哪次采集”。

PS 使用 \nolinkurl{DualMaskFullFrameTracker} 在滤波掩膜上提取轮廓、选取目标，并计算面积、凸包、凸缺陷和边界框。通过缺陷夹角、深度及形状指标识别握拳、剪刀和张开等状态；结合选中轮廓内的原始掩膜统计得到目标质心。分类稳定性处理再生成 hand\_valid、位置、手势类别和事件序列。[E05]

视觉输入按“连续状态”处理：采集侧逐帧核验身份，识别侧消费最新待处理帧，避免业务负载变化时累积历史位置。掩膜复制到 CPU 自有不可变内存后交付，DMA 缓冲可以按拥有权协议继续流转。位置用于方向与悬停；新的握拳事件用于确认；保持握拳用于画布连续创作。状态与事件的分离使菜单、游戏和画布获得一致的操作语义。[E06、E13]

\subsection{帧边界启停状态机}
视觉采集由 \nolinkurl{vision_source_admission} 控制源接入。该模块在视频时钟域实现四状态 FSM，将软件的打开、关闭请求转换为以帧边界为单位的动作。同步后的关闭请求优先处理；RUNNING 与 CLOSING 均保持 axis\_enable，使停止请求可以等待当前帧完整结束。[E20]

\begin{figure}[htbp]\centering\begin{tikzpicture}[node distance=23mm and 47mm]
\node[box,text width=34mm] (wait) {WAIT\_START\\等待帧边界};
\node[box,text width=34mm,right=of wait] (run) {RUNNING\\接纳图像};
\node[box,text width=34mm,below=of run] (closing) {CLOSING\\等待当前帧结束};
\node[box,text width=34mm,below=of wait] (closed) {CLOSED\\源端关闭};
\draw[arr] (wait)--node[lab,above]{open + 有效边界}(run);
\draw[arr] (run)--node[lab,right]{close\\尚未到边界}(closing);
\draw[arr] (closing)--node[lab,below]{有效边界}(closed);
\draw[arr] (wait)--node[lab,left]{close}(closed);
\draw[arr] (run.south west)--node[lab,sloped,above]{close + 有效边界}(closed.north east);
\draw[ctl] (closed.west)--++(-9mm,0)|-node[lab,near start,left]{复位}(wait.west);
\node[font=\footnotesize,below=8mm of closed.south west,anchor=north west,text width=130mm] {有效边界：vid\_vsync 上升沿且 vid\_active\_video = 0。未满足转移条件时保持原状态。};
\end{tikzpicture}
\caption{采集源四状态机}\label{fig:3}\end{figure}
{\small\renewcommand{\arraystretch}{1.28}
\begin{longtable}{@{}>{\raggedright\arraybackslash}p{0.24\linewidth}>{\raggedright\arraybackslash}p{0.35\linewidth}>{\raggedright\arraybackslash}p{0.34\linewidth}@{}}
\toprule
\bfseries 当前状态 & \bfseries 转移条件 & \bfseries 下一状态与作用\\\midrule
\endfirsthead
\toprule
\bfseries 当前状态 & \bfseries 转移条件 & \bfseries 下一状态与作用\\\midrule
\endhead
\bottomrule\endfoot
WAIT\_START & 打开请求且到达有效帧边界 & RUNNING；开始接纳图像流\\
WAIT\_START & 关闭请求 & CLOSED；完成尚未开始的取消\\
RUNNING & 关闭请求且当前为有效帧边界 & CLOSED；在边界直接关闭\\
RUNNING & 关闭请求且尚未到边界 & CLOSING；继续接纳当前帧\\
CLOSING & 到达有效帧边界 & CLOSED；停止后续帧接入\\
CLOSED & 复位 & WAIT\_START；为下一次采集初始化\\
\end{longtable}}
有效帧边界定义为 vid\_vsync 上升沿且 \nolinkurl{vid_active_video} 为低。若边界与有效像素同时出现，逻辑记录 boundary\_fault。关闭后还需等待延迟的像素与帧计数事件结清，\nolinkurl{source_quiescent} 才成立。采集生命周期进一步结合流水线排空、DMA 排空及缓冲拥有者状态，形成从源端到内存端的完整停止确认。

\subsection{三槽缓冲与帧拥有权}
\nolinkurl{vision_frame_owner} 以三个循环槽连接硬件生产与 CPU 消费。设累计授予帧数为 G，累计释放帧数为 R，正常运行保持 0 $\leq$ G-R $\leq$ 3。仅当 G-R < 3 时允许授予新槽，从而防止生产者覆盖仍被 CPU 使用的帧。[E21]

CPU 可领取一帧的条件是：存在已授予未释放帧、当前未持有其他槽，且 RGB、掩膜包、状态三路完成计数均已越过当前消费位置。领取和释放使用 epoch 与 frame\_id 组成的令牌；释放成功后推进消费者槽与 R。完成计数负责数据就绪，令牌负责身份匹配，二者共同维护同一帧的生命周期。

{\small\renewcommand{\arraystretch}{1.28}
\begin{longtable}{@{}>{\raggedright\arraybackslash}p{0.24\linewidth}>{\raggedright\arraybackslash}p{0.35\linewidth}>{\raggedright\arraybackslash}p{0.34\linewidth}@{}}
\toprule
\bfseries 阶段 & \bfseries 生产者/消费者行为 & \bfseries 硬件保证\\\midrule
\endfirsthead
\toprule
\bfseries 阶段 & \bfseries 生产者/消费者行为 & \bfseries 硬件保证\\\midrule
\endhead
\bottomrule\endfoot
授予 & 硬件取得空槽并开始本帧写入 & 在途帧数最多为 3\\
就绪 & RGB、掩膜、状态三路完成 & 完成前不向 CPU 开放领取\\
领取 & CPU 提交本帧令牌，持有消费槽 & 同一时刻最多一个 CPU 持有槽\\
释放 & CPU 完成复制并提交匹配令牌 & 释放后槽才可再次分配\\
排空 & 源帧数固定，各完成计数与授予/释放数一致 & 与 DMA 排空联合确认生命周期结束\\
\end{longtable}}
RTL 对重复领取、重复释放、令牌不符、计数倒退和配置中途改变设置粘滞故障。相应 testbench 包含连续七帧、CPU 持有 64 周期、三路完成、令牌检查和计数边界等场景。设计把缓冲生命周期落实为可验证的状态与计数，而不是依赖软件对时间先后的估计。[E24]

\subsection{Scene2D 图形处理与 HDMI 输出}
\begin{figure}[htbp]\centering\begin{tikzpicture}[node distance=10mm and 10mm]
\node[soft] (scene) {PS 应用状态\\Scene 编译};
\node[box,right=of scene] (dma) {DDR + AXI DMA\\命令 / 资源};
\node[box,right=of dma] (render) {PL 命令执行\\裁剪 / 光栅化};
\node[box,below=12mm of render] (back) {后台缓冲\\320×180 token};
\node[box,left=of back] (front) {前台缓冲\\320×180 token};
\node[box,left=of front] (scan) {调色板 + 4× 放大\\HDMI 720p60};
\draw[arr] (scene)--(dma);\draw[arr] (dma)--(render);\draw[arr] (render)--(back);
\draw[arr,<->] (back)--node[lab,above=5mm]{帧边界交换}(front);\draw[arr] (front)--(scan);
\node[soft,below=12mm of front,text width=73mm] (done) {提交序号与完成计数\\接收 → 渲染 → 呈现};
\draw[ctl] (back.south)|-(done.east);\draw[ctl] (scan.south)|-(done.west);
\end{tikzpicture}
\caption{命令式图形处理与双缓冲呈现}\label{fig:4}\end{figure}
Scene2D 是 PL 中的命令式二维绘图处理器。PS 将应用状态编译成清屏、图块层、索引色贴图、矩形和结束帧命令；AXI DMA 从 DDR 取出命令包，PL 完成解码、裁剪、资源读取和像素写入。菜单、四游戏、画布与演示页面共用这一绘制接口。[E07–E08]

逻辑分辨率为 320×180。源图像使用 4 bit 调色板索引，索引 0 表示透明；帧缓冲采用 8 bit token，高 4 bit 选择调色板组，低 4 bit 选择组内颜色。显示端查表得到 RGB888，再按 4 倍最近邻放大到 1280×720。双逻辑帧缓冲的有效像素容量为 115,200 B。

渲染写入后台缓冲，像素扫描读取前台缓冲。跨时钟握手将就绪帧交给显示域，在垂直消隐起点交换前后台身份。RGB、同步信号和有效信号保持匹配的流水延迟，输出端以 74.25 MHz 像素时钟生成 60 Hz HDMI 时序。扫描读取片上存储，使显示节奏独立于 PS 业务调度。[E25]

软件为每次提交分配 submit\_seq，并依次核对包接收、渲染完成与呈现完成。render\_cycles 提供渲染周期，render\_status 提供执行状态，\nolinkurl{present_completed_seq/count} 表示显示端已接纳目标帧。这些寄存器使图形接口具备从命令到显示的明确完成语义。

静态背景、字体与模板可驻留复用。页面悬停只追加四条边框矩形命令；应用更新普通场景时传输的是紧凑命令，而无需每次重新组织整幅 RGB 图像。内容更新保留在 PS，像素处理保留在 PL，兼顾了功能扩展与固定时序输出。

\subsection{已完成的应用功能}
{\small\renewcommand{\arraystretch}{1.28}
\begin{longtable}{@{}>{\raggedright\arraybackslash}p{0.24\linewidth}>{\raggedright\arraybackslash}p{0.35\linewidth}>{\raggedright\arraybackslash}p{0.34\linewidth}@{}}
\toprule
\bfseries 功能 & \bfseries 输入与状态变化 & \bfseries 输出与存储约束\\\midrule
\endfirsthead
\toprule
\bfseries 功能 & \bfseries 输入与状态变化 & \bfseries 输出与存储约束\\\midrule
\endhead
\bottomrule\endfoot
主菜单与内容列表 & 手位置悬停，新握拳确认，返回控件停留 & 共享场景主题；\nolinkurl{HOME/GAMES/VIDEOS} 使用同一 live 手拥有者\\
四游戏 & 统一位置、手势、事件与按键输入 & 公共 GPU 输出；Game1 为 BTN0 盾、BTN1 攻击\\
视频 & 列表选择与播放/退出状态 & 两段默认视频；播放前完成机械手释放交接\\
手势演示 & 进入示范页后按演示序列运行 & 0–9 与爱心的实体手/像素手对应展示；示范独占\\
握拳画布 & 位置、握拳、形状及尺寸选择 & 图形合并到固定 46,080 字节画布，不积累对象列表\\
临时相册 & 保存、查看、返回 & 只保留最近一次保存；完整初始化或重启后清空\\
暂停与结果页 & 公共返回/暂停与结果选择 & 全屏替换页面，统一后续流转\\
\end{longtable}}
四游戏共享 GameInput、GamePlugin 生命周期与场景编译器。统一输入既能来自板端真实设备，也能用于 PC/PYNQ 相同序列的逐 tick 对照。Game1 使用 BTN0 触发盾、BTN1 触发攻击；菜单用悬停与新握拳确认，返回采用公共停留控件。[E01、E13]

画布将绘制结果合并到固定 46,080 B 的像素存储中，持续创作保持有界内存。落笔采用与界面一致的握拳判定；限频时保留最新输入，在真正处理后推进消费身份。相册保存最近一次作品，完整初始化或重启后清空。公共加载、暂停和结果页统一各应用的进入、退出及继续流程。

\subsection{灵巧手协同与模式管理}
\begin{figure}[htbp]\centering\begin{tikzpicture}[node distance=10mm and 10mm]
\node[soft] (obs) {视觉 / 游戏观察\\来源与序列};
\node[box,right=of obs] (owner) {独立手进程\\动作拥有者};
\node[box,right=of owner] (motor) {ORCA / USB 串口\\17 个电机};
\node[soft,below=14mm of obs] (live) {菜单 / 列表\\live 跟随};
\node[soft,right=of live] (pause) {画布 / 相册\\暂停动作};
\node[soft,right=of pause] (demo) {示范模式\\独占姿态序列};
\node[box,below=10mm of pause,text width=65mm] (off) {视频 / 预览交接\\17OFF 读回 → CLOSED 回执};
\draw[arr] (obs)--node[lab,above=6mm]{本机 UDP}(owner);\draw[arr] (owner)--node[lab,above=6mm]{1 Mbaud}(motor);
\draw[ctl] (owner)--(pause);\draw[ctl] (owner.south west)--(live.north east);
\draw[ctl] (owner.south east)--(demo.north west);\draw[ctl] (pause)--(off);
\end{tikzpicture}
\caption{灵巧手观察链路与模式交接}\label{fig:5}\end{figure}
灵巧手由观察发布器、同机数据报接口、独立解释器进程、动作拥有者、ORCA 适配层与 USB 串口链路组成。实体端使用 scs225、1 Mbaud 串口与 ID 1–17 的电机。视觉发布器给出最新手势快照，动作进程完成来源检查、姿态编译、标定映射和目标下发。[E10–E11]

HOME、GAMES、VIDEOS 列表共享 live 拥有者；画布与相册暂停动作；示范独占播放 0–9 与爱心的实体手/像素手对应内容；视频与预览交接先完成全部 17 个电机 OFF。返回 live 时重新建立新鲜观察来源。动作执行期间保存最新目标，当前阶段结束后选择最新观察，使动作链跟随当前交互意图。

使能前读取电机状态，检查位置范围、电压、温度与扭矩状态，并设置速度和当前目标；执行时检查通信返回，跳过未变化的目标；退出时逐个写入扭矩关闭并读回确认。v2.0 的串口验证已覆盖 18 项协议检查及两次 17OFF/CLOSED 会话回执，形成可进入、可运行、可退出的设备生命周期。[E02]

\clearpage
\section{软硬件划分依据与接口设计}
\subsection{划分依据与模块职责}
划分围绕四个判断展开：数据规模、计算规律、修改频率和实时性。逐像素重复运算与扫描时序放在 PL；不规则轮廓处理、手势语义和应用状态放在 PS；内容编辑与测试观测放在 PC。系统选择双掩膜、场景命令、手势观察作为稳定的模块边界，使各端通过紧凑协议协作。

{\small\renewcommand{\arraystretch}{1.28}
\begin{longtable}{@{}>{\raggedright\arraybackslash}p{0.2\linewidth}>{\raggedright\arraybackslash}p{0.15\linewidth}>{\raggedright\arraybackslash}p{0.29\linewidth}>{\raggedright\arraybackslash}p{0.27\linewidth}@{}}
\toprule
\bfseries 模块 & \bfseries 位置 & \bfseries 选择依据 & \bfseries 配套设计要求\\\midrule
\endfirsthead
\toprule
\bfseries 模块 & \bfseries 位置 & \bfseries 选择依据 & \bfseries 配套设计要求\\\midrule
\endhead
\bottomrule\endfoot
摄像头接入与预处理 & PL & 连续图像流；相同局部操作可流水化 & 必须处理时钟跨越、帧边界和源故障\\
双掩膜封装及 DMA & PL + PS 驱动 & 批量数据搬运与同帧身份检查 & 包格式、缓冲拥有权与缓存一致性必须共同约定\\
轮廓与手势判断 & PS & 轮廓数与长度不固定；规则需快速调整 & 受 CPU 负载影响；最新帧策略抑制积压\\
游戏、菜单、画布 & PS & 状态分支多，内容变化频繁 & 受 SDK、命令容量和逻辑 tick 约束\\
命令光栅化与扫描 & PL & 大量重复绘制；像素输出需固定时序 & 受 BRAM、命令集和调色板格式约束\\
手势 OSD & PS 发布 + PL 叠加 & 输入提示独立于游戏场景变化 & 要求新鲜度、提交握手和像素坐标一致\\
按键 & PL 同步 + PS 消抖 & 固定外部引脚与灵活事件策略 & 轮询受调度影响；不能等同硬实时中断\\
灵巧手 & PS 独立进程 & 复用串口 SDK、标定及生命周期管理 & 需严格单拥有者和释放确认\\
Studio & PC & 编辑、打包、对照和状态观测 & 网络断开不等于板端资源已释放\\
\end{longtable}}
视觉前端将局部像素操作流水化，减少 PS 对每个像素重复执行的工作；Scene2D 将规则绘图与扫描留在硬件，PS 专注页面和游戏决策。轮廓长度与数量具有变化性，保留在 Python 侧便于调整规则。灵巧手独立进程则复用既有串口与标定依赖，并集中管理唯一动作来源。

这一分工支持作品持续扩展：增加游戏或调整菜单主要修改 PS 插件与资源；改变图像阈值通过寄存器配置；绘图和视频时序由同一 Overlay 提供。公共 SDK 连接内容开发和硬件能力，使多个应用共享一套已验证底座。

\subsection{总线结构与寄存器映射}
PS 通过 GP0 访问 AXI-Lite 配置及状态外设。场景和资源上传由 axi\_dma 的 MM2S 经 HP0 读取 DDR；RGB 捕获由 frame\_dma 的 S2MM 经 HP1 写入 DDR；双掩膜通过 moments\_dma 的 S2MM/SG 经 HP3 回写 DDR。批量数据与控制寄存器的分离使 CPU 可以用少量读写管理整帧事务。[E14]

{\small\renewcommand{\arraystretch}{1.28}
\begin{longtable}{@{}>{\raggedright\arraybackslash}p{0.24\linewidth}>{\raggedright\arraybackslash}p{0.35\linewidth}>{\raggedright\arraybackslash}p{0.34\linewidth}@{}}
\toprule
\bfseries 实例及基地址 & \bfseries 本版本角色 & \bfseries 关键约束\\\midrule
\endfirsthead
\toprule
\bfseries 实例及基地址 & \bfseries 本版本角色 & \bfseries 关键约束\\\midrule
\endhead
\bottomrule\endfoot
axi\_dma 0x40400000 & 场景与图形资源上传 & MM2S；提交完成与 DMA 完成分开检查\\
moments\_dma 0x40420000 & 双掩膜回传 & S2MM/SG；包身份和 CRC 验证\\
build\_id\_gpio 0x41200000 & 硬件身份及按键 & +0x00 读 BUILD\_ID；+0x08 的 bit0/1 读 BTN0/1\\
\nolinkurl{packet_status_gpio} 0x41220000 & 场景包状态 & 上传接受/拒绝计数检查\\
\nolinkurl{render_completion_gpio} 0x41270000 & 渲染完成 & +0x00 序号；+0x08 完成计数\\
\nolinkurl{render_cycles_gpio} 0x41290000 & 渲染周期 & 驱动读取当次完成周期\\
\nolinkurl{render_status_gpio} 0x412A0000 & 渲染状态 & 正常提交要求状态为 0\\
\nolinkurl{present_completion_gpio} 0x412B0000 & 显示完成 & +0x00 序号；+0x08 呈现计数\\
capture\_control 0x413F0000 & 采集上下文与生命周期控制 & 身份、参数回读、ARM/中止/屏障\\
capture\_status 0x41410000 & 采集 epoch 与源故障 & 必须一致读取状态及 epoch\\
camera\_iic 0x41600000 & 摄像头配置 & 使用配对的配置及驱动\\
frame\_dma 0x43000000 & RGB 捕获通路 & 主要用于需要 RGB 的路径；与掩膜消费分离\\
demosaic 0x43C00000 & 去马赛克控制 & 配置匹配视频尺寸与链路\\
input\_osd\_0 0x43C20000 & 独立手势标记 & 4 KiB 地址窗口；忙时不覆盖提交\\
\end{longtable}}
驱动根据配对的 bit/HWH、运行清单与 BUILD\_ID 建立 Overlay 会话，视觉、OSD 和图形模块在该会话中借用硬件。R23b2 将 BTN0/BTN1 接入 BUILD\_ID GPIO 的第二输入通道，复用已有 AXI 从设备，使按键扩展保持简洁。[E15]

\subsection{AXI 流握手与背压设计}
AXI-Stream 在时钟上升沿同时满足 TVALID 与 TREADY 时传输一拍，记 take = TVALID $\land$ TREADY。发送端在 TVALID=1、TREADY=0 时保持数据及附属标志；接收端在空间不足时拉低 TREADY。视觉流水线用这一规则连接局部运算、FIFO 和输出配对，保持数据与身份一致。[E19、E24]

{\small\renewcommand{\arraystretch}{1.28}
\begin{longtable}{@{}>{\raggedright\arraybackslash}p{0.2\linewidth}>{\raggedright\arraybackslash}p{0.14\linewidth}>{\raggedright\arraybackslash}p{0.14\linewidth}>{\raggedright\arraybackslash}p{0.43\linewidth}@{}}
\toprule
\bfseries 时钟拍示例 & \bfseries TVALID & \bfseries TREADY & \bfseries 数据与动作\\\midrule
\endfirsthead
\toprule
\bfseries 时钟拍示例 & \bfseries TVALID & \bfseries TREADY & \bfseries 数据与动作\\\midrule
\endhead
\bottomrule\endfoot
k & 1 & 1 & A 被接收；发送端可推进到 B\\
k+1 & 1 & 0 & B 保持；像素计数不推进\\
k+2 & 1 & 0 & B 与其帧标志继续保持\\
k+3 & 1 & 1 & B 被接收；数据和帧标志同时推进\\
k+4 & 0 & 1 & 无有效数据；不发生传输\\
\end{longtable}}
双掩膜配对要求原始与滤波两路同时有效，且下游可以接纳新结果。上下文 FIFO 在首个被接收的输入像素入队，在最终被接收的输出帧尾出队。本文所述双掩膜内部帧尾 tlast 按该模块定义使用；视频接入接口的行、帧标志由各自适配模块转换。

AXI-Lite 则承担离散寄存器事务。软件先写完整参数，再写提交寄存器，并通过 busy、状态或回读核对结果。OSD 使用“参数—提交—忙闲”结构，采集使用“参数—回读—ARM—完成/排空”结构，避免半更新配置被数据通路读取。

\subsection{双掩膜 DMS1 数据接口}
DMS1 为小端格式。每个像素依次携带原始掩膜与滤波掩膜两个字节，随后追加 16 个 uint32 组成的 64 B 尾部。640×360 的像素负载为 460,800 B，完整包为 460,864 B；一般长度为 ceil(width×height/2)×4+64，奇数像素按零填充对齐。[E04]

{\small\renewcommand{\arraystretch}{1.28}
\begin{longtable}{@{}>{\raggedright\arraybackslash}p{0.24\linewidth}>{\raggedright\arraybackslash}p{0.35\linewidth}>{\raggedright\arraybackslash}p{0.34\linewidth}@{}}
\toprule
\bfseries 尾部字序号 & \bfseries 内容 & \bfseries 校验目的\\\midrule
\endfirsthead
\toprule
\bfseries 尾部字序号 & \bfseries 内容 & \bfseries 校验目的\\\midrule
\endhead
\bottomrule\endfoot
W0–W1 & 0x31534D44（DMS1）、0x00100001 & 包类型、版本与结构长度\\
W2–W4 & frame\_id、config\_id、epoch & 帧身份、阈值配置及采集会话\\
W5 & height 在高 16 bit；width 在低 16 bit & 几何尺寸与解码形状一致\\
W6–W8 & 像素数、负载字节数、状态 0 & 防止截断或异常结果被消费\\
W9–W10 & 六个 HSV 阈值字节 & 与本次采集配置完全一致\\
W11 & PAIR\_FORMAT = 0x00010802 & 确认原始/滤波掩膜配对格式\\
W12–W14 & 保留，必须为 0 & 拒绝未支持的格式扩展\\
W15 & CRC32 & 覆盖之前所有包字节，不含 CRC 字本身\\
\end{longtable}}
软件联合校验长度、CRC、几何、阈值、frame\_id、config\_id 和 epoch。原始掩膜为 0/255 二值，滤波结果允许灰度值。通过后生成 CPU 自有不可变快照，交给识别线程。该格式把两种图像结果和配置身份封装成一次完整交付。

采集控制上下文写入 capture\_control 的 +0x20、+0x24、+0x28、+0x2C、+0x30，内容包含六个 HSV 阈值、帧号、配置号、epoch 及对齐字段。软件回读参数后发出 ARM，使用非零 epoch 标识本次会话。capture\_status 的 +0x00 为 0x31545356，+0x04 为 ABI 0x00010000，+0x08 为 \nolinkurl{ready/sealed/fault}，+0x0C 为原因码，+0x10 为 epoch；读取状态前后核对 epoch，以形成一致状态快照。[E16]

\subsection{Scene2D 场景、资源与完成接口}
场景包采用小端 uint32，固定为 64 B 头和 4,096 B 命令区，总计 4,160 B。每条命令 8 字、32 B，最多 128 条，command\_count 指定有效条数。[E07–E08]

{\small\renewcommand{\arraystretch}{1.28}
\begin{longtable}{@{}>{\raggedright\arraybackslash}p{0.29\linewidth}>{\raggedright\arraybackslash}p{0.66\linewidth}@{}}
\toprule
\bfseries 场景头字序号 & \bfseries 字段\\\midrule
\endfirsthead
\toprule
\bfseries 场景头字序号 & \bfseries 字段\\\midrule
\endhead
\bottomrule\endfoot
W0–W1 & S2DG 魔数；ABI major/minor = 1.0\\
W2–W4 & frame\_id、submit\_seq、command\_count\\
W5–W7 & camera\_x、camera\_y、RGB888 背景色\\
W8–W11 & \nolinkurl{palette_generation}、\nolinkurl{predicted_cycles}、\nolinkurl{asset_generation}、\nolinkurl{tilemap_generation}\\
W12–W15 & 保留，填 0\\
\end{longtable}}
{\small\renewcommand{\arraystretch}{1.28}
\begin{longtable}{@{}>{\raggedright\arraybackslash}p{0.24\linewidth}>{\raggedright\arraybackslash}p{0.35\linewidth}>{\raggedright\arraybackslash}p{0.34\linewidth}@{}}
\toprule
\bfseries 命令 & \bfseries 编码 & \bfseries 执行含义\\\midrule
\endfirsthead
\toprule
\bfseries 命令 & \bfseries 编码 & \bfseries 执行含义\\\midrule
\endhead
\bottomrule\endfoot
CLEAR & 1 & 清除逻辑画布并建立背景\\
DRAW\_TILE\_LAYER & 2 & 从 tilemap 与图块资源绘制层\\
BLIT\_INDEX4 & 3 & 索引色图像拷贝，支持定义的透明与翻转语义\\
FILL\_RECT & 4 & 在裁剪范围内写入矩形 token\\
END\_FRAME & 255 & 结束本次命令序列\\
\end{longtable}}
资源包由 64 B 头与实际负载组成，采用 \nolinkurl{S2DA/S2DT/S2DP/S2DF} 魔数区分类别。W1 为 ABI，W2 为类型，W3 为 slot，W4 为 generation，W5 为负载字节数，W6 为负载 CRC32，W7 为 flags，W8–W15 保留。资源包 CRC 覆盖负载，DMS1 CRC 则覆盖 CRC 字之前的整个包，两者按各自协议计算。

资源代数 generation 将命令与其所需图集、调色板和 tilemap 对应起来。PS 完成资源上传并等待代数可见，再提交使用该代数的 Scene。图集 bundle 使用三个 16 B 描述符，分别描述 sprite、tile、font，共享数据容量为 32,768 B。FRAME8 接口提供 57,600 B 完整 token 帧传输，普通游戏场景采用命令式绘制。

提交按接收、渲染、呈现三个阶段确认。接收计数说明包已被硬件接受；渲染序号和状态说明命令已执行；呈现序号说明显示端已切换到该结果。超时或状态异常时，驱动结束当前 sink 的使用并按会话清理流程恢复，使后续提交保持清晰的序号关系。

\subsection{独立手势 OSD 与板载按键}
OSD 将输入反馈从应用场景中独立出来：PS 按 30 Hz 目标发布周期提供位置与握拳状态，PL 将标记叠加到 60 Hz HDMI 输出。识别视图为 640×360，写入寄存器前映射为 1280×720 像素坐标。[E12]

{\small\renewcommand{\arraystretch}{1.28}
\begin{longtable}{@{}>{\raggedright\arraybackslash}p{0.24\linewidth}>{\raggedright\arraybackslash}p{0.35\linewidth}>{\raggedright\arraybackslash}p{0.34\linewidth}@{}}
\toprule
\bfseries OSD 偏移 & \bfseries 内容 & \bfseries 提交规则\\\midrule
\endfirsthead
\toprule
\bfseries OSD 偏移 & \bfseries 内容 & \bfseries 提交规则\\\midrule
\endhead
\bottomrule\endfoot
+0x00 & 身份 0x4F534401 & 打开服务时校验\\
+0x04 & busy 状态 bit0 & 忙时跳过本次，不覆盖在途数据\\
+0x08 & 提交寄存器 & 先写完整参数，再写 1 提交\\
+0x10 & bit0 enabled；bit1 valid；bit2 fist & valid/fist 受新鲜度控制\\
+0x14 & x 低 16 bit；y 高 16 bit & HDMI 输出坐标\\
+0x18 & 0 & 兼容字段置零\\
+0x1C & TTL 放在高 16 bit & 按 60 Hz 帧期换算并限制上界\\
+0x20 & 观察序列 sequence & 区分新观察与重复发布\\
+0x24 & 0 & 保留/历史字段置零\\
\end{longtable}}
观察年龄小于 250 ms 时参与有效性判断，TTL 按剩余时间换算为显示帧并限制至最多 15。忙时跳过本次发布并统计，待空闲后提交最新观察。这样，即使应用场景没有变化，用户仍能看到独立更新的手势标记。

BTN0/BTN1 位于 build\_id\_gpio +0x08 的低两位。PL 同步外部引脚，PS 按 4 ms 目标周期采样，连续三个相同采样确认稳定状态，仅在上升沿产生带 sequence 的 ButtonSignal。按键事件进入容量 64 的有界队列，连续位置则采用最新值策略，分别适配离散事件和连续状态。[E17]

\subsection{GamePlugin 与 Studio 网络协议}
GamePlugin 通过 reset(context)、step(game\_input)、snapshot() 组织生命周期。GameInput 提供 tick\_id、timestamp\_ns、hand\_valid、[-1,1] 范围的 hand\_x/hand\_y、握拳事件与序列、手势类别与变化序列，以及 actions 和 button。持续状态描述当前输入，事件序列保证一次确认或技能只消费一次。[E13]

Studio 与板端 Agent 通过 TCP 通信，默认端口 8765。VVGP 1.0 使用大端序 48 B 头，布局为 >4sHHHHIIQQIII，最大负载 8 MiB。[E09]

{\small\renewcommand{\arraystretch}{1.28}
\begin{longtable}{@{}>{\raggedright\arraybackslash}p{0.24\linewidth}>{\raggedright\arraybackslash}p{0.35\linewidth}>{\raggedright\arraybackslash}p{0.34\linewidth}@{}}
\toprule
\bfseries 头偏移及长度 & \bfseries 字段 & \bfseries 作用\\\midrule
\endfirsthead
\toprule
\bfseries 头偏移及长度 & \bfseries 字段 & \bfseries 作用\\\midrule
\endhead
\bottomrule\endfoot
0 / 4 B & VVGP & 协议魔数\\
4、6 / 各 2 B & major、minor & 双端版本匹配\\
8、10 / 各 2 B & message\_type、flags & 消息类别与标志\\
12、16 / 各 4 B & header\_size、payload\_size & 分帧及长度限制\\
20、28 / 各 8 B & sequence、timestamp\_ns & 消息身份及时间信息\\
36、40 / 各 4 B & source\_id、channel\_id & 来源与通道\\
44 / 4 B & checksum & 校验字段置零后的头部与负载 CRC32\\
\end{longtable}}
协议包含 HELLO/HELLO\_ACK、\nolinkurl{COMMAND/COMMAND_RESULT}、SUBSCRIBE、DATA、TELEMETRY、PING/PONG 与 ERROR。接收器在 TCP 字节流中累积完整消息，再按长度和 CRC 解码；握手交换 SDK 与输入 schema 能力。Studio 因而能够统一组织部署、订阅、测试和状态回传，板端继续承担应用实际运行。

\subsection{灵巧手观察、动作与释放接口}
主进程通过 127.0.0.1 UDP 向独立手进程发送 JSON，schema 为 hand-passive-v1，最大包长 2,048 B。外层包含 schema、session、sequence、revoked、value；value 携带视觉或游戏来源与观察。发送采用非阻塞模式，接收端检查会话、递增序列、来源与默认 0.5 s 新鲜度。[E10]

run\_id 与 config\_epoch 标识视觉来源，frame\_id 标识新观察。ORCA 适配层将姿态映射为电机位置，检查 0–4095 编码范围，并通过 SCS225 客户端下发目标和速度。冻结参数 speed=500 按驱动单位使用。释放时向电机寄存器 40 写 0 并读回，汇总各 ID 状态。[E11]

{\small\renewcommand{\arraystretch}{1.28}
\begin{longtable}{@{}>{\raggedright\arraybackslash}p{0.24\linewidth}>{\raggedright\arraybackslash}p{0.35\linewidth}>{\raggedright\arraybackslash}p{0.34\linewidth}@{}}
\toprule
\bfseries 阶段 & \bfseries 接口回执 & \bfseries 系统用途\\\midrule
\endfirsthead
\toprule
\bfseries 阶段 & \bfseries 接口回执 & \bfseries 系统用途\\\midrule
\endhead
\bottomrule\endfoot
观察接收 & 格式、会话、序列与新鲜度校验 & 选择本次有效输入\\
姿态编译 & 生成通过范围检查的电机目标 & 把手势语义转换为设备参数\\
动作下发 & 串口 SDK 通信返回 & 确认命令传输过程\\
阶段结束 & 当前动作流程结束，选择最新目标 & 衔接后续动作\\
17OFF & ID 1–17 扭矩关闭写入和读回 & 完成设备释放检查\\
CLOSED & 会话关闭回执 & 完成业务模式交接\\
\end{longtable}}
独立解释器隔离 PYNQ 图像链路与 ORCA 的依赖，窄观察协议连接两端。业务模式切换通过同一动作拥有者协调，因此菜单跟随、创作暂停、演示独占与视频释放具有统一的控制入口。

\subsection{时钟、复位与跨域约束}
工程按采集、控制/渲染与显示划分时钟域。源接入把打开、关闭请求同步到视频域，图形端把帧就绪与切换确认在渲染域和像素域之间握手传递；OSD 使用邮箱结构提交完整参数。各通路在跨域处明确数据保持、请求同步及完成确认条件。[E20、E22–E23]

{\small\renewcommand{\arraystretch}{1.28}
\begin{longtable}{@{}>{\raggedright\arraybackslash}p{0.24\linewidth}>{\raggedright\arraybackslash}p{0.35\linewidth}>{\raggedright\arraybackslash}p{0.34\linewidth}@{}}
\toprule
\bfseries 时钟/信号域 & \bfseries 实现报告中的约束 & \bfseries 对应职责\\\midrule
\endfirsthead
\toprule
\bfseries 时钟/信号域 & \bfseries 实现报告中的约束 & \bfseries 对应职责\\\midrule
\endhead
\bottomrule\endfoot
clk\_fpga\_0 & 100 MHz；10 ns & AXI 控制与主数据处理/渲染\\
像素时钟 & 74.250 MHz；13.468 ns & 720p 扫描\\
SerialClkIO & 371.250 MHz；2.694 ns & HDMI 串行化相关时钟\\
clk\_fpga\_1 & 23.810 MHz；41.999 ns & PS 输出的另一时钟\\
vid\_pclk & 100 MHz；10 ns 的输入约束 & 摄像头输入时钟分析\\
\end{longtable}}
源启停请求使用三级同步器，寄存器标记 ASYNC\_REG 并禁止移位寄存器抽取；视频复位通过 XPM 复位同步结构处理。这样，异步控制先进入稳定的本地时钟状态，再参与帧边界 FSM。

OSD 邮箱跨域采用两类约束配合：请求与确认同步器入口设置定点 false\_path；64 bit 数据邮箱设置 10 ns 的 max\_delay -datapath\_only 和 10 ns 的 bus\_skew。载荷在握手期间保持稳定，同步后的请求通知接收端取用完整数据。约束脚本检查源、目的端点数量，确保约束落实到预期寄存器，避免层次变化后静默失效。[E22]

板级 XDC 将 BTN0/BTN1 约束到 D19/D20，I/O 标准为 LVCMOS33；摄像头 IIC 使用 U9/T9 与上拉。vid\_pclk 网络采用 \nolinkurl{CLOCK_DEDICATED_ROUTE} FALSE 处理板级时钟接入布线需求，时序仍由相应时钟与路径约束分析。

R23b2 路由后报告给出 WNS +0.011 ns、TNS 0、WHS +0.050 ns、THS 0，setup/hold 违例端点均为 0。该结果对应本次发布的实际硬件基线，表明已施加约束下的建立与保持检查收敛。完整报告同时保留外部 I/O 延迟覆盖检查，相关说明集中列于附录。[E23]

\subsection{完整交互事务与容量设计}
“进入游戏并释放技能”贯穿了各层接口：双掩膜包通过身份与 CRC 校验；PS 识别手位置并更新菜单悬停；新握拳触发进入游戏与资源初始化；插件在 tick 内更新状态并生成 Scene；DMA 上传后由硬件渲染并给出呈现回执；BTN0 的消抖上升沿形成独立序列事件，Game1 更新盾技能状态并生成下一帧反馈。

“画布落笔并保存”复用同一视觉与显示链路：界面握拳判定成立，创作页在频率预算内消费最新观察，将图形合并到固定画布；形状和大小选择通过悬停边框反馈；保存后生成最近一次相册内容。两个例子说明，稳定接口可以组合出不同的应用体验，而硬件核心保持复用。

{\small\renewcommand{\arraystretch}{1.28}
\begin{longtable}{@{}>{\raggedright\arraybackslash}p{0.24\linewidth}>{\raggedright\arraybackslash}p{0.35\linewidth}>{\raggedright\arraybackslash}p{0.34\linewidth}@{}}
\toprule
\bfseries 项目 & \bfseries 按接口定义计算 & \bfseries 设计用途\\\midrule
\endfirsthead
\toprule
\bfseries 项目 & \bfseries 按接口定义计算 & \bfseries 设计用途\\\midrule
\endhead
\bottomrule\endfoot
单逻辑帧 token & 320×180×1 = 57,600 B & 片上像素有效容量\\
双帧缓冲 & 2×57,600 = 115,200 B & 不含调色板、命令、图集及 BRAM 对齐开销\\
固定场景包 & 64 + 128×32 = 4,160 B & 不含额外资源上传\\
若每秒提交 60 场景 & 4,160×60 = 249,600 B/s & 命令传输预算\\
640×360 双掩膜包 & 640×360×2 + 64 = 460,864 B & 包含校验尾部\\
若双掩膜为 30 帧/s & 460,864×30 = 13,825,920 B/s & 给定帧率条件下的带宽预算\\
100 MHz 的 60 Hz 帧期 & 100,000,000/60 $\approx$ 1,666,667 周期 & 绘制与控制共用的时钟周期预算\\
\end{longtable}}
表中容量由已定义接口直接计算，带帧率条件的行用于系统预算。双缓冲使绘制和扫描并行，命令式上传减小普通场景的主机传输量，最新帧策略控制输入积压，有界画布控制长期内存增长。这些设计分别对应显示连续性、数据搬运、输入响应和创作运行四个实际需求。

\subsection{分层验证与作品完成度}
v2.0 采用“RTL 验证点—软件合同—真实设备—完整业务流程”的分层验证方式。底层 testbench 针对握手保持、帧身份与缓冲拥有权设置场景；应用以同输入逐 tick 比对验证 PC/PYNQ 一致性；真实 GPU 和串口检查则验证设备路径；最后通过服务与 Studio 重连确认运行入口。[E02、E24]

{\small\renewcommand{\arraystretch}{1.28}
\begin{longtable}{@{}>{\raggedright\arraybackslash}p{0.24\linewidth}>{\raggedright\arraybackslash}p{0.35\linewidth}>{\raggedright\arraybackslash}p{0.34\linewidth}@{}}
\toprule
\bfseries 层次 & \bfseries 已有成果 & \bfseries 与系统设计的对应关系\\\midrule
\endfirsthead
\toprule
\bfseries 层次 & \bfseries 已有成果 & \bfseries 与系统设计的对应关系\\\midrule
\endhead
\bottomrule\endfoot
RTL 验证设计 & 流背压/配置上下文测试；frame\_owner 14 类场景 & 数据、标志、令牌与缓冲生命周期\\
四游戏对照 & 16 条 PC/PYNQ 用例 PASS & 统一输入合同与插件状态一致性\\
图形与创作 & 真实 GPU 60 项检查 PASS & 创作、示范、暂停/结果和重启路径\\
真实串口接口 & 18 项协议检查 PASS & 菜单跟随与业务模式交接\\
设备释放 & 两次会话完成 17OFF/CLOSED & 动作进程的关闭与资源回收\\
服务运行 & 默认服务、测试注入移除、Studio 重连通过 & 日常运行与开发管理入口\\
硬件实现 & R23b2 建立/保持检查无违例 & 当前发布 Overlay 的实现收敛\\
\end{longtable}}
作品已形成从 FPGA 逻辑到用户功能的连续实现：双掩膜提供统一视觉入口，Scene2D 提供统一图形出口，GamePlugin 连接四游戏及创作场景，独立动作通道连接实体手，Studio 连接开发与验证。功能扩展、状态管理、设备交接和回执检查都建立在同一 v2.0 基线上。

\clearpage
\section{附录：版本、证据与测试口径}
\subsection{发布身份与数据分类}
发布提交为 \nolinkurl{bd6858bbe4ded19caf1f91322850150c97cab21e}，运行清单 SHA-256 为 \nolinkurl{5f53d85d5bcdb8c871d387683cbf7b6b12c899b380b9370642c9bd2ff41e9d56}。正文中的格式、尺寸与容量来自源码；WNS/WHS 来自实现报告；16/60/18 与关闭回执来自 v2.0 验收。RTL testbench 的场景列表来自冻结测试源文件。

v2.0 正式技能输入为板载按钮；兼容代码中的 Voice View 名称与 tone/volume 字段保留历史接口。实体手本轮证据为协议与 SDK 验证，\nolinkurl{acceptance-summary} 中 \nolinkurl{mechanical_verified} 为 false；真人姿态测量、摄像头到屏幕统一时延及全新 SD 安装属于独立验证项目。本文交互闭环描述用户操作和系统反馈，电机位置误差控制由执行器层处理。

时序报告 check\_timing 记录 4 个 no\_input\_delay 和 6 个 no\_output\_delay。正文的建立/保持结果针对已约束路径；外部 I/O 时序覆盖应结合完整约束清单阅读。vid\_pclk 表项为输入时钟约束值。带宽预算、OSD 发布目标与消抖采样周期分别属于设计值，不作为识别吞吐或整机实测时延。

冻结候选 manifest 中的 release\_ready=false、board\_pending=true 保留生成时的阶段记录；本报告功能验收采用其后的 v2.0 \nolinkurl{acceptance-summary}。视觉运行时注明改编自 \nolinkurl{KPTKleist/Pynq_Gesture_Detect}（MIT），灵巧手适配使用 ORCA 代码与配置；相应版权与许可随工程保留。

\subsection{来源定位}
下表以 B 表示 \nolinkurl{build/v2-release-repo/releases/v2.0/board_runtime/}，S 表示同版本 source/，V 表示 verification/，P 表示 \nolinkurl{B/python/pynq_multimodal_game/}。完整路径、文件大小与 SHA-256 见随附 \nolinkurl{evidence_index.json}。

{\small\renewcommand{\arraystretch}{1.28}
\begin{longtable}{@{}>{\raggedright\arraybackslash}p{0.29\linewidth}>{\raggedright\arraybackslash}p{0.66\linewidth}@{}}
\toprule
\bfseries 编号 & \bfseries 文件定位与用途\\\midrule
\endfirsthead
\toprule
\bfseries 编号 & \bfseries 文件定位与用途\\\midrule
\endhead
\bottomrule\endfoot
E01 & \nolinkurl{tools/release_v2/RELEASE_NOTES.md}；正式功能及修复范围\\
E02 & \nolinkurl{V/acceptance-summary.json}；16/60/18、关闭及验收边界\\
E03 & \nolinkurl{docs/bringup/2026-10-29_v2_release.md}；封存与发布身份\\
E04 & \nolinkurl{B/board/pynq_z2/dual_mask_packet.py}；DMS1 格式与校验\\
E05 & \nolinkurl{B/board/pynq_z2/scripts/vision_async_runtime.py}；轮廓识别与双掩膜消费\\
E06 & \nolinkurl{B/board/pynq_z2/vision_latest_capture.py}；最新帧与来源检查\\
E07 & P/gpu\_abi.py、P/scene.py；图形接口与画面格式\\
E08 & \nolinkurl{P/agent/c6i_video.py}、\nolinkurl{P/agent/c6i_backend.py}；提交与会话\\
E09 & P/protocol.py、\nolinkurl{P/agent/transport.py}；VVGP 与 TCP\\
E10 & \nolinkurl{P/hand_shadow_bridge.py}、\nolinkurl{P/hand_serial_session.py}；\nolinkurl{B/hand_child/serial_motion_receiver.py}；\nolinkurl{B/hand-runtime/src/pynq_usb_hand/game_bridge.py}\\
E11 & \nolinkurl{B/hand-runtime/src/pynq_usb_hand/orca_backend.py}；\nolinkurl{B/hand-runtime/assets/config.yaml}；动作适配与电机配置\\
E12 & P/button\_osd.py；位置、TTL 与寄存器提交\\
E13 & P/game\_sdk.py；\nolinkurl{P/creative_pages.py}；插件、输入与创作\\
E14 & \nolinkurl{B/overlay/c6v_candidate.hwh}、\nolinkurl{board_test_manifest.json}；真实地址与硬件签核\\
E15 & \nolinkurl{B/docs/HANDOFF.md}；R22 至 R23b2 的结构取舍记录\\
E16 & \nolinkurl{B/board/pynq_z2/vision_capture_control.py}；采集控制事务\\
E17 & \nolinkurl{P/button_input.py}；按键消抖、序列与队列\\
E18 & \nolinkurl{docs/scene2d_gpu_architecture.md}；基础图形架构，具体 ABI 以 E07/E08 为准\\
E19 & \nolinkurl{S/rtl/vision_streaming/axis_rgb_mask_frontend.sv}、\nolinkurl{axis_binary_mask_pair.sv}；视觉数据通路\\
E20 & \nolinkurl{S/rtl/vision_streaming/vision_source_admission.sv}、\nolinkurl{vision_capture_lifecycle.sv}；启停 FSM\\
E21 & \nolinkurl{S/rtl/vision_streaming/vision_frame_owner.sv}；三槽拥有权\\
E22 & \nolinkurl{S/subprojects/button_only/constraints/button_only.xdc}、\nolinkurl{button_osd_cdc.tcl}；引脚与 CDC 约束\\
E23 & \nolinkurl{B/overlay/timing.rpt}；R23b2 路由后时序\\
E24 & \nolinkurl{S/tests/rtl/tb_vision_frame_owner.sv}、\nolinkurl{tb_axis_rgb_mask_frontend.sv}；测试场景定义\\
E25 & \nolinkurl{S/rtl/scene2d_gpu/scene2d_c4m7_core.sv}、\nolinkurl{scene2d_c4m3_core.sv}、\nolinkurl{scene2d_frame_swap.sv}；图形核心与帧交换\\
\end{longtable}}
\subsection{官方指南与总报告衔接}
R01 为用户指定的官方仓库自主选题指南，读取日期 2026-10-30。原文保存在 \nolinkurl{sources/open_topic_guide.md}，SHA-256 为 \nolinkurl{bd57e4f8b911bd7b872eab51000576512b542410a7eca6166c31d319de0d5678}。原文地址见下方链接。指南注明，完整指南与大赛官方公告具有最终解释优先级。

\url{https://gitee.com/Vickyiii/fpgachina26-amd/blob/master/tracks/open-topic/GUIDE.md}
本报告对应总报告的系统设计两章；其余章节衔接创新点、优化对比、AI 协作、Skill 与复现说明，工程提交附开源许可及目录对照。


\end{document}
